% GENERATED FILE. Edit canonical lessons, then run npm run generate:books.
% canonical-source-hash: fnv1a64:e963ba3db59b4455
% canonical-lessons: MR-C237-parampara, MR-C237-boli, MR-C237-prasiddhi, MR-C237-sanketsthal, MR-C237-kalphalak

\chapter{Tradition, Dialect, Publicity, Website, Keyboard}
\label{ch:mr-tradition-dialect-publicity-website-keyboard}
\hlchaptermodality{237}

\begin{chapteropening}
\textbf{By the end of this chapter:} \emph{I can say a tradition, a dialect, publicity, a website and a keyboard.}

Complete the last lesson of chapter 237: I can say a tradition, a dialect, publicity, a website and a keyboard.
\end{chapteropening}

\section[\texorpdfstring{\mr{परंपरा} (paramparā)}{परंपरा (paramparā)}]{\mr{परंपरा} (paramparā) — a tradition}

\label{lesson:MR-C237-parampara}



\begin{quote}
Before the new one: say the Marathi for a booking, then the Marathi for a river valley.
\end{quote}

\subsection*{You'll want to know: \mr{परंपरा}}
\textbf{\mr{परंपरा}} — \emph{paramparā} — \textquotedblleft{}a tradition\textquotedblright{}.

\begin{grammarlens}[title={What you've built}]
One more word for this chapter.
\end{grammarlens}

\subsection*{Your turn}
\begin{itemize}
\raggedright
  \item \emph{Say it:} \emph{paramparā}
  \item \emph{Say it:} \emph{paramparā}, once more
  \item \emph{Say it:} say \emph{paramparā}
  \item \emph{Recall:} say the Marathi for a booking, then the Marathi for a river valley, then say \emph{paramparā} again
\end{itemize}

\begin{tcolorbox}[breakable,colback=teal!4,colframe=teal!35!black,title={Before you move on}]
What does \mr{परंपरा} mean? (\textquotedblleft{}A tradition\textquotedblright{}.) Say it once more.
\end{tcolorbox}
\section[\texorpdfstring{\mr{बोली} (bolī)}{बोली (bolī)}]{\mr{बोली} (bolī) — a dialect}

\label{lesson:MR-C237-boli}



\begin{quote}
Before the new one: say the Marathi for a river valley, then the Marathi for a tradition.
\end{quote}

\subsection*{You'll want to know: \mr{बोली}}
\textbf{\mr{बोली}} — \emph{bolī} — \textquotedblleft{}a dialect\textquotedblright{}.

\begin{grammarlens}[title={What you've built}]
One more word for this chapter.
\end{grammarlens}

\subsection*{Your turn}
\begin{itemize}
\raggedright
  \item \emph{Say it:} \emph{bolī}
  \item \emph{Say it:} \emph{bolī}, once more
  \item \emph{Say it:} say \emph{bolī}
  \item \emph{Recall:} say the Marathi for a river valley, then the Marathi for a tradition, then say \emph{bolī} again
\end{itemize}

\begin{tcolorbox}[breakable,colback=teal!4,colframe=teal!35!black,title={Before you move on}]
What does \mr{बोली} mean? (\textquotedblleft{}A dialect\textquotedblright{}.) Say it once more.
\end{tcolorbox}
\section[\texorpdfstring{\mr{प्रसिद्धी} (prasiddhī)}{प्रसिद्धी (prasiddhī)}]{\mr{प्रसिद्धी} (prasiddhī) — publicity, fame}

\label{lesson:MR-C237-prasiddhi}



\begin{quote}
Before the new one: say the Marathi for a tradition, then the Marathi for a dialect.
\end{quote}

\subsection*{You'll want to know: \mr{प्रसिद्धी}}
\textbf{\mr{प्रसिद्धी}} — \emph{prasiddhī} — \textquotedblleft{}publicity, fame\textquotedblright{}.

\begin{grammarlens}[title={What you've built}]
One more word for this chapter.
\end{grammarlens}

\subsection*{Your turn}
\begin{itemize}
\raggedright
  \item \emph{Say it:} \emph{prasiddhī}
  \item \emph{Say it:} \emph{prasiddhī}, once more
  \item \emph{Say it:} say \emph{prasiddhī}
  \item \emph{Recall:} say the Marathi for a tradition, then the Marathi for a dialect, then say \emph{prasiddhī} again
\end{itemize}

\begin{tcolorbox}[breakable,colback=teal!4,colframe=teal!35!black,title={Before you move on}]
What does \mr{प्रसिद्धी} mean? (\textquotedblleft{}Publicity, fame\textquotedblright{}.) Say it once more.
\end{tcolorbox}
\section[\texorpdfstring{\mr{संकेतस्थळ} (saṅketsthaḷ)}{संकेतस्थळ (saṅketsthaḷ)}]{\mr{संकेतस्थळ} (saṅketsthaḷ) — a website}

\label{lesson:MR-C237-sanketsthal}



\begin{quote}
Before the new one: say the Marathi for a dialect, then the Marathi for publicity.
\end{quote}

\subsection*{You'll want to know: \mr{संकेतस्थळ}}
\textbf{\mr{संकेतस्थळ}} — \emph{saṅketsthaḷ} — \textquotedblleft{}a website\textquotedblright{}.

\begin{grammarlens}[title={What you've built}]
One more word for this chapter.
\end{grammarlens}

\subsection*{Your turn}
\begin{itemize}
\raggedright
  \item \emph{Say it:} \emph{saṅketsthaḷ}
  \item \emph{Say it:} \emph{saṅketsthaḷ}, once more
  \item \emph{Say it:} say \emph{saṅketsthaḷ}
  \item \emph{Recall:} say the Marathi for a dialect, then the Marathi for publicity, then say \emph{saṅketsthaḷ} again
\end{itemize}

\begin{tcolorbox}[breakable,colback=teal!4,colframe=teal!35!black,title={Before you move on}]
What does \mr{संकेतस्थळ} mean? (\textquotedblleft{}A website\textquotedblright{}.) Say it once more.
\end{tcolorbox}
\section[\texorpdfstring{\mr{कळफलक} (kaḷphalak)}{कळफलक (kaḷphalak)}]{\mr{कळफलक} (kaḷphalak) — a keyboard}

\label{lesson:MR-C237-kalphalak}



\begin{quote}
Before the new one: say the Marathi for publicity, then the Marathi for a website.
\end{quote}

\subsection*{You'll want to know: \mr{कळफलक}}
\textbf{\mr{कळफलक}} — \emph{kaḷphalak} — \textquotedblleft{}a keyboard\textquotedblright{}.

\begin{grammarlens}[title={What you've built}]
One more word for this chapter.
\end{grammarlens}

\subsection*{Your turn}
\begin{itemize}
\raggedright
  \item \emph{Say it:} \emph{kaḷphalak}
  \item \emph{Say it:} \emph{kaḷphalak}, once more
  \item \emph{Say it:} say \emph{kaḷphalak}
  \item \emph{Recall:} say the Marathi for publicity, then the Marathi for a website, then say \emph{kaḷphalak} again
\end{itemize}

\begin{tcolorbox}[breakable,colback=teal!4,colframe=teal!35!black,title={Before you move on}]
What does \mr{कळफलक} mean? (\textquotedblleft{}A keyboard\textquotedblright{}.) Say it once more.
\end{tcolorbox}
